% from ega4-16-fr.tex
\oldpage[IV-4]{24}
Notons d'abord que $g_{X/Y/Z}$ correspond au cas où dans (16.4.3.3) on fait
$X'=X$, $S'=Y$, $S=Z$ et $u=1_X$, et $f_{X/Y/Z}$ au cas où l'on remplace
$X'$, $X$, $S$, $S'$ par $X$, $Y$, $Z$, $Z$ respectivement et $u$, $f$ par
$f$, $g$ respectivement.

On a un diagramme commutatif (I, 5.3.5)
\[
  \xymatrix{
    X\ar[r]^-{\Delta_f}\ar[dr]_-f
      &X\times_Y X\ar[d]^-p\ar[r]^-j
      &X\times_Z X\ar[d]^-{f\times_Z f}\\
    &Y\ar[r]_-{\Delta_g}&Y\times_Z Y
  }
  \tag{16.4.18.3}
  \label{IV.16.4.18.3}
\]
où $j=(1_X,1_X)_Z$ est une immersion, $j\circ\Delta_f=\Delta_{g\circ f}$ et
$p$ est le morphisme structural. Comme on peut se borner au cas où $X$, $Y$,
$Z$ sont affines, on peut supposer les immersions $\Delta_f$, $\Delta_g$ et
$j$ fermées, de sorte que $\mathcal{O}_X$ et
$\mathcal{O}_{X\times_Y X}$ s'identifient respectivement à
$\mathcal{O}_{X\times_Z X}/\mathscr{J}$ et
$\mathcal{O}_{X\times_Z X}/\mathscr{L}$, où
$\mathscr{J}\supset\mathscr{L}$ sont les deux Idéaux quasi-cohérents
correspondant respectivement aux immersions $\Delta_{g\circ f}$ et $j$. La
$\mathcal{O}_X$-Algèbre $\mathscr{P}_{X/Z}^n$ s'identifie donc à
$\mathcal{O}_{X\times_Z X}/\mathscr{J}^{n+1}$, et
$\mathscr{P}_{X/Y}^n$ s'identifie à
$\mathcal{O}_{X\times_Y X}/(\mathscr{J}/\mathscr{L})^{n+1}$, c'est-à-dire à
$\mathcal{O}_{X\times_Z X}/(\mathscr{J}^{n+1}+\mathscr{L})$, et par suite au
quotient de $\mathscr{P}_{X/Z}^n$ par
$(\mathscr{J}^{n+1}+\mathscr{L})/\mathscr{J}^{n+1}$. Mais on sait
(\emph{loc. cit.}) que $p$ et $j$ font de $X\times_Y X$ le produit des
$(Y\times_Z Y)$-préschémas $Y$ et $X\times_Z X$, donc si $\mathcal{O}_Y$ est
identifié à $\mathcal{O}_{Y\times_Z Y}/\mathscr{K}$, où $\mathscr{K}$ est
l'Idéal correspondant à $\Delta_g$, $\mathscr{L}$ est égal à
$(f\times_Z f)^*(\mathscr{K})\cdot\mathcal{O}_{X\times_Z X}$ (I, 4.4.5).
Comme $(\mathscr{J}^{n+1}+\mathscr{L})/\mathscr{J}^{n+1}$ est l'Idéal de
$\mathscr{P}_{X/Z}^n$ engendré par l'image de $\mathscr{L}$, on en déduit la
proposition.

\begin{corollary}[16.4.19]
\label{IV.16.4.19}
Avec les notations de (16.4.18), on a une suite exacte de
$\mathcal{O}_X$-Modules quasi-cohérents
\[
  \xymatrix{
    f^*(\Omega_{Y/Z}^1)\ar[r]^-{f_{X/Y/Z}}&
    \Omega_{X/Z}^1\ar[r]^-{g_{X/Y/Z}}&
    \Omega_{X/Y}^1\ar[r]&0.
  }
  \tag{16.4.19.1}
  \label{IV.16.4.19.1}
\]
\end{corollary}

Lorsque $X$, $Y$, $Z$ sont affines, on retrouve ainsi la suite exacte
(0, 20.5.7.1).

\begin{proposition}[16.4.20]
\label{IV.16.4.20}
Soient $f:Y\to Z$ un morphisme, $j:X\to Y$ une immersion fermée,
$\mathscr{K}$ l'Idéal quasi-cohérent de $\mathcal{O}_Y$ correspondant à $j$.
Alors on a
$\mathscr{P}_{X/Y}^n=\mathcal{O}_X=\mathcal{O}_Y/\mathscr{K}$,
l'homomorphisme canonique
$j_{X/Y/Z}:j^*(\mathscr{P}_{Y/Z}^n)\to\mathscr{P}_{X/Z}^n$ est surjectif, et
son noyau est l'Idéal de $j^*(\mathscr{P}_{Y/Z}^n)$ engendré par
$j^*(\mathcal{O}_Y\cdot d_{Y/Z}^n(\mathscr{K}))$ (il faut noter que
$d_{Y/Z}^n(\mathscr{K})$ est un sous-faisceau de groupes commutatifs de
$\mathscr{P}_{Y/Z}^n$, mais non un $\mathcal{O}_Y$-Module en général).
\end{proposition}

On sait (I, 5.3.8) que la diagonale
$\Delta_j:X\to X\times_Y X$ est un isomorphisme, d'où la première assertion.
Si $\varpi_1$ et $\varpi_2$ sont les deux homomorphismes d'Algèbres
$\mathcal{O}_Y\to\mathscr{P}_{Y/Z}^n$ correspondant respectivement aux deux
projections canoniques $p_1$, $p_2$ de $Y\times_Z Y$ sur $Y$, rappelons que
par définition (16.3.5 et 16.3.6) $\varpi_1$ est l'homomorphisme structural de
la $\mathcal{O}_Y$-Algèbre $\mathscr{P}_{Y/Z}^n$ et
$\varpi_2=d_{Y/Z}^n$. La $\mathcal{O}_X$-Algèbre
$j^*(\mathscr{P}_{Y/Z}^n)$ s'identifie donc à
$\mathscr{P}_{Y/Z}^n/\varpi_1(\mathscr{K})\mathscr{P}_{Y/Z}^n$ et son quotient
par l'Idéal engendré par $j^*(d_{Y/Z}^n(\mathscr{K}))$ à
\[
  \mathscr{P}_{Y/Z}^n/
  (\varpi_1(\mathscr{K})+\varpi_2(\mathscr{K}))\mathscr{P}_{Y/Z}^n.
\]
Notons maintenant qu'on a un diagramme commutatif
